% GENERATED FILE. Edit canonical lessons, then run npm run generate:books.
% canonical-source-hash: fnv1a64:459df64dcc0c5a2b
% canonical-lessons: PA-C217-came-return, PA-C217-aaia, PA-C217-aaia-review, PA-C217-aurat, PA-C217-au, PA-C217-who-came, PA-C217-payoff

\chapter{Who Came?}
\label{ch:pa-a2-aaia}
\hlchaptermodality{217}

\begin{chapteropening}
\textbf{By the end of this chapter:} \emph{I can say who came, a man, a woman or a group, and tell a man from a woman in the past.}

Say who came with aia, ai or ae, use aurat 'a woman', and recognise the letter au.
\end{chapteropening}

\section[\texorpdfstring{\pa{ਆਉਣਾ} · \pa{ਮੈਨੂੰ} \pa{ਸਮਝ} … (āuṇā · mainūṁ …)}{ਆਉਣਾ · ਮੈਨੂੰ ਸਮਝ … (āuṇā · mainūṁ …)}]{Return to come}

\label{lesson:PA-C217-came-return}



\begin{quote}
\emph{Hear:} \emph{āuṇā} Say \textbf{\pa{ਆਉਣਾ}}, \textquotedblleft{}to come\textquotedblright{}. Then \textbf{\pa{ਮੈਨੂੰ} \pa{ਸਮਝ} \pa{ਨਹੀਂ} \pa{ਆਈ}}.
\end{quote}

\begin{grammarlens}[title={Grammar Lens}]
Listen to the last word of \textbf{\pa{ਮੈਨੂੰ} \pa{ਸਮਝ} \pa{ਨਹੀਂ} \pa{ਆਈ}}: \emph{āī}. It is a came form, and you have been saying it since the early lessons. Next to it, hold the went forms \textbf{\pa{ਗਿਆ} · \pa{ਗਈ}}.
\end{grammarlens}

\subsection*{Your turn}
\begin{itemize}
\raggedright
  \item \emph{Hear:} \emph{to come} Say \textbf{\pa{ਆਉਣਾ}}.
  \item \emph{Hear:} \emph{Aman went home} Say \textbf{\pa{ਅਮਨ} \pa{ਘਰ} \pa{ਗਿਆ}}.
  \item \emph{Hear:} \emph{I did not understand} Say \textbf{\pa{ਮੈਨੂੰ} \pa{ਸਮਝ} \pa{ਨਹੀਂ} \pa{ਆਈ}}.
\end{itemize}

\begin{tcolorbox}[breakable,colback=teal!4,colframe=teal!35!black,title={Before you move on}]
Say \emph{went}, for a man and a woman, once more.
\end{tcolorbox}
\section[\texorpdfstring{\pa{ਆਇਆ} · \pa{ਆਈ} · \pa{ਆਏ} (āiā · āī · āe)}{ਆਇਆ · ਆਈ · ਆਏ (āiā · āī · āe)}]{\pa{ਆਇਆ}, \pa{ਆਈ}, \pa{ਆਏ} — came}

\label{lesson:PA-C217-aaia}



\begin{quote}
Say \emph{to come}. Then \emph{Aman went home yesterday}.
\end{quote}

\begin{grammarlens}[title={Grammar Lens: came, three ways}]
\textbf{\pa{ਆਉਣਾ}} has three came forms, like the three went forms: \textbf{\pa{ਆਇਆ}} for a man, \textbf{\pa{ਆਈ}} for a woman, \textbf{\pa{ਆਏ}} for a group or for respect.

\begin{quote}
\textbf{\pa{ਅਮਨ} \pa{ਕੱਲ੍ਹ} \pa{ਆਇਆ।}} — \emph{Aman kallh āiā.} — Aman came yesterday.
\end{quote}

\begin{quote}
\textbf{\pa{ਕੁੜੀ} \pa{ਕੱਲ੍ਹ} \pa{ਆਈ।}} — \emph{kuṛī kallh āī.} — The girl came yesterday.
\end{quote}

\begin{quote}
\textbf{\pa{ਉਹ} \pa{ਕੱਲ੍ਹ} \pa{ਆਏ।}} — \emph{uh kallh āe.} — They came yesterday.
\end{quote}

The \textbf{\pa{ਆਈ}} in \textbf{\pa{ਮੈਨੂੰ} \pa{ਸਮਝ} \pa{ਨਹੀਂ} \pa{ਆਈ}} is the woman's form: \textbf{\pa{ਸਮਝ}}, \textquotedblleft{}understanding\textquotedblright{}, is a feminine word, and it is the one that did not come. CDIAL lists Punjabi \emph{āiā}, \textquotedblleft{}he came\textquotedblright{} (CDIAL 1045).
\end{grammarlens}

\subsection*{Your turn}
\begin{itemize}
\raggedright
  \item \emph{Hear:} \emph{the boy came yesterday} Say \textbf{\pa{ਮੁੰਡਾ} \pa{ਕੱਲ੍ਹ} \pa{ਆਇਆ।}}
  \item \emph{Hear:} \emph{they came home} Say \textbf{\pa{ਉਹ} \pa{ਘਰ} \pa{ਆਏ।}}
  \item \emph{Read it:} \textbf{\pa{ਕੁੜੀ} \pa{ਕੱਲ੍ਹ} \pa{ਆਈ।}} Give its meaning.
  \item \emph{Write it:} copy \textbf{\pa{ਆਇਆ}} once
\end{itemize}

\begin{tcolorbox}[breakable,colback=teal!4,colframe=teal!35!black,title={Before you move on}]
\emph{Came}, for a man? (\textbf{\pa{ਆਇਆ}}.) A woman? (\textbf{\pa{ਆਈ}}.) A group? (\textbf{\pa{ਆਏ}}.)
\end{tcolorbox}
\section[\texorpdfstring{\pa{ਆਇਆ} · \pa{ਆਈ} · \pa{ਆਏ} (āiā · āī · āe)}{ਆਇਆ · ਆਈ · ਆਏ (āiā · āī · āe)}]{Hear and say came again}

\label{lesson:PA-C217-aaia-review}



\begin{quote}
\emph{Hear:} \emph{āiā; āī; āe} Say \textbf{\pa{ਆਇਆ} · \pa{ਆਈ} · \pa{ਆਏ}}.
\end{quote}

\subsection*{Your turn}
\begin{itemize}
\raggedright
  \item \emph{Hear:} \emph{Aman came yesterday} Say \textbf{\pa{ਅਮਨ} \pa{ਕੱਲ੍ਹ} \pa{ਆਇਆ।}}
  \item \emph{Hear:} \emph{the girl went home} Say \textbf{\pa{ਕੁੜੀ} \pa{ਘਰ} \pa{ਗਈ।}}
  \item \emph{Read it:} \textbf{\pa{ਉਹ} \pa{ਕੱਲ੍ਹ} \pa{ਆਏ।}} Give its meaning.
  \item \emph{Write it:} cover the model and write \textbf{\pa{ਆਈ}} once
\end{itemize}

\begin{tcolorbox}[breakable,colback=teal!4,colframe=teal!35!black,title={Before you move on}]
Say who came yesterday, and who went.
\end{tcolorbox}
\section[\texorpdfstring{\pa{ਔਰਤ} (aurat)}{ਔਰਤ (aurat)}]{\pa{ਔਰਤ} (aurat) — a woman}

\label{lesson:PA-C217-aurat}



\begin{quote}
Before the new word: say \emph{a man}, \textbf{\pa{ਆਦਮੀ}}. Then \emph{the girl came yesterday}.
\end{quote}

\subsection*{You'll want to know: \pa{ਔਰਤ}}
\textbf{\pa{ਔਰਤ}} — \emph{aurat} — \textquotedblleft{}a woman\textquotedblright{}. A feminine noun, so the past that agrees with her ends like \textbf{\pa{ਆਈ}}:

\begin{quote}
\textbf{\pa{ਔਰਤ} \pa{ਕੱਲ੍ਹ} \pa{ਆਈ।}} — \emph{aurat kallh āī.} — The woman came yesterday.
\end{quote}

\begin{cousinweb}
\textbf{\pa{ਔਰਤ}} (\emph{aurat}) is a loan: it is usually traced to Persian \emph{aurat}, \textquotedblleft{}a woman\textquotedblright{}, a word Persian itself took from Arabic. CDIAL gives no Punjabi form.
\end{cousinweb}

\begin{grammarlens}[title={What you've built}]
One more word for talking about the past.
\end{grammarlens}

\subsection*{Your turn}
\begin{itemize}
\raggedright
  \item \emph{Say it:} \emph{aurat}
  \item \emph{Hear:} \emph{the woman came yesterday} Say \textbf{\pa{ਔਰਤ} \pa{ਕੱਲ੍ਹ} \pa{ਆਈ।}}
  \item \emph{Recall:} say \emph{a man}, then \emph{a woman}
  \item \emph{Write it:} copy \textbf{\pa{ਔਰਤ}} once from the model
\end{itemize}

\begin{tcolorbox}[breakable,colback=teal!4,colframe=teal!35!black,title={Before you move on}]
What does \textbf{\pa{ਔਰਤ}} mean? (\textquotedblleft{}A woman\textquotedblright{}.) Say it once more.
\end{tcolorbox}
\section[\texorpdfstring{\pa{ਔ} (au)}{ਔ (au)}]{\pa{ਔ} — one vowel from the word for a woman}

\label{lesson:PA-C217-au}



\begin{quote}
Say \textbf{\pa{ਔਰਤ}}, \textquotedblleft{}a woman\textquotedblright{}, the word from the previous lesson.
\end{quote}

\subsection*{Script you'll notice: \pa{ਔ}}
\begin{quote}
\textbf{\pa{ਅ}} + \textbf{\pa{ੌ}} $\to$ \textbf{\pa{ਔ}}
\end{quote}

\begin{quote}
\textbf{\pa{ਔ}} — \emph{au} — an independent Gurmukhi vowel
\end{quote}

You know \textbf{\pa{ਅ}} and the sign \textbf{\pa{ੌ}}. Together they make the standing vowel \textbf{\pa{ਔ}}, the first letter of \textbf{\pa{ਔਰਤ}}.

No stroke order is taught here: the track does not yet have a verified source for one.

\subsection*{Writing: copy what you can see}
Keep \textbf{\pa{ਔ}} in front of you. Trace its printed shape once with a finger, then copy it once with a pencil while the model stays visible. Say \emph{au} when you finish.

\begin{quote}
This is observe-and-copy with the model visible, not writing from memory.
\end{quote}

\subsection*{Your turn}
\begin{itemize}
\raggedright
  \item \emph{Look:} at \textbf{\pa{ਔ}}; YOU SAY: \emph{au}
  \item \emph{Read it:} \textbf{\pa{ਔਰਤ}}, then point to \textbf{\pa{ਔ}}
  \item \emph{Write it:} copy \textbf{\pa{ਔ}} once while the model remains visible
\end{itemize}

\begin{tcolorbox}[breakable,colback=teal!4,colframe=teal!35!black,title={Before you move on}]
Point to \textbf{\pa{ਔ}}, say \emph{au}, then say the whole word \textbf{\pa{ਔਰਤ}}.

Source: \href{https://www.unicode.org/charts/PDF/U0A00.pdf}{Unicode Gurmukhi chart}.
\end{tcolorbox}
\section[\texorpdfstring{\pa{ਕੌਣ} \pa{ਆਇਆ}? (kauṇ āiā?)}{ਕੌਣ ਆਇਆ? (kauṇ āiā?)}]{\pa{ਕੌਣ} \pa{ਆਇਆ}? — who came?}

\label{lesson:PA-C217-who-came}



\begin{quote}
Say \emph{who?}, \textbf{\pa{ਕੌਣ}}. Then \emph{the woman came yesterday}.
\end{quote}

\begin{grammarlens}[title={Grammar Lens: the one who came decides}]
Asking, you do not know who came, so you use the plain \textbf{\pa{ਆਇਆ}}. The answer fits the person:

\begin{quote}
\textbf{\pa{ਕੌਣ} \pa{ਆਇਆ}?} — \emph{kauṇ āiā?} — Who came?
\end{quote}

\begin{quote}
\textbf{\pa{ਆਦਮੀ} \pa{ਆਇਆ।}} — \emph{ādmī āiā.} — A man came.
\end{quote}

\begin{quote}
\textbf{\pa{ਔਰਤ} \pa{ਆਈ।}} — \emph{aurat āī.} — A woman came.
\end{quote}

The went forms work the same way: \textbf{\pa{ਆਦਮੀ} \pa{ਗਿਆ}}, \textbf{\pa{ਔਰਤ} \pa{ਗਈ}}.
\end{grammarlens}

\subsection*{Your turn}
\begin{itemize}
\raggedright
  \item \emph{Hear:} \emph{a woman went home} Say \textbf{\pa{ਔਰਤ} \pa{ਘਰ} \pa{ਗਈ।}}
  \item \emph{Hear:} \emph{a man came yesterday} Say \textbf{\pa{ਆਦਮੀ} \pa{ਕੱਲ੍ਹ} \pa{ਆਇਆ।}}
  \item \emph{Read it:} \textbf{\pa{ਔਰਤ} \pa{ਆਈ।}} Point to \textbf{\pa{ਔ}}, then give the meaning.
  \item \emph{Write it:} copy \textbf{\pa{ਔਰਤ} \pa{ਆਈ।}} once
\end{itemize}

\begin{tcolorbox}[breakable,colback=teal!4,colframe=teal!35!black,title={Before you move on}]
Ask \emph{who came?} Then answer: a woman came.
\end{tcolorbox}
\section[\texorpdfstring{\pa{ਕੌਣ} \pa{ਆਇਆ}? (kauṇ āiā?)}{ਕੌਣ ਆਇਆ? (kauṇ āiā?)}]{Who came: a four-skill payoff}

\label{lesson:PA-C217-payoff}



\begin{quote}
Say \emph{came} three ways. Then \emph{a woman}.
\end{quote}

\subsection*{Your turn: four skills}
\begin{itemize}
\raggedright
  \item \emph{Hear:} \emph{who came yesterday?} Say \textbf{\pa{ਕੱਲ੍ਹ} \pa{ਕੌਣ} \pa{ਆਇਆ}?}
  \item \emph{Hear:} \emph{a woman came} Say \textbf{\pa{ਔਰਤ} \pa{ਆਈ।}}
  \item \emph{Read it:} \textbf{\pa{ਆਦਮੀ} \pa{ਘਰ} \pa{ਗਿਆ।}} Give its meaning, then point to \textbf{\pa{ਔ}} in \textbf{\pa{ਔਰਤ}}.
  \item \emph{Write it:} cover the line and write \textbf{\pa{ਔਰਤ} \pa{ਕੱਲ੍ਹ} \pa{ਆਈ।}}
\end{itemize}

\begin{tcolorbox}[breakable,colback=teal!4,colframe=teal!35!black,title={Before you move on}]
Tell who came to your home yesterday, a man, a woman or a group.
\end{tcolorbox}
